\section{Graph pullback, the complete wall drift, and a finite core repair}
\label{sec:r13-graph-pullback-wall-drift}

A continuous piecewise-affine coordinate graph can repair a
coefficient obstruction by changing the full incident spaces.
Pulling its primitive back to the old real branch generally
creates normal jumps.  This section records their exact measure,
quantifies the drift required by PK, and separates a universal
real graph repair from the extra source branches of actual units.

\subsection{The full distributional graph chain rule}

Let \(U\subset\mathbf R^d\) be open and let
\(F:U\to\mathbf R^m\) be continuous and affine on each
cell of a finite polyhedral decomposition.  Let
\(B=\sum_{a=1}^m B_a(y)\,dy_a\) be a smooth covector on a
neighborhood of \(F(U)\).  On cell interiors define
\[
 g=F^*B=(DF)^TB(F),\qquad
 A=(DF)^T e(B)(F)\,DF.
 \tag{GD1}
\]
The cellwise field \(g\) is locally bounded and belongs to
\(\mathrm{BD}_{\rm loc}(U)\).  For an internal facet \(H\),
choose the unit normal \(\nu_H\) pointing from its minus
cell to its plus cell.  Continuity of \(F\) gives
\[
             DF^+-DF^-=c_H\otimes\nu_H,
                       \qquad c_H\in\mathbf R^m.
 \tag{GD2}
\]
A common refinement includes every internal facet of every
coordinate, including facets whose individual contributions
subsequently cancel.

\begin{proposition}[The weighted wall measure]
\label{prop:r13-graph-full-strain-chain-rule}
The full distributional symmetric derivative is
\[
 \begin{split}
 Eg&=A\,\mathcal L^d+\mathcal W_{F,B},\\
 \mathcal W_{F,B}
    &=\sum_{a=1}^m B_a(F)\,D^2F_a\\
    &=\sum_H\langle B(F),c_H\rangle\,
             \nu_H\otimes\nu_H\,
                     \mathcal H^{d-1}\!\llcorner H.
 \end{split}
 \tag{GD3}
\]
Multiplication in the middle line is by the continuous
function \(B_a(F)\).  All formulas remain valid when
\(DF\) has deficient rank.
\end{proposition}

\begin{proof}
On a cell the ordinary chain rule gives the regular symmetric
derivative \(A\).  The one-sided traces of \(g\) differ by
\[
             g^+-g^-=\langle B(F),c_H\rangle\nu_H.
 \tag{GD4}
\]
The distributional derivative of a piecewise smooth field is
its regular derivative plus its jump tensor with the oriented
facet normal.  Taking the symmetric part of GD4 gives the
last line of GD3.  There are no extra contributions on
intersections of two or more facets: the first derivative of
a locally bounded piecewise smooth field has the displayed
facet measure, and those intersections have zero facet
measure.  Applying the same jump calculation to the
piecewise constant gradient of each scalar \(F_a\) gives
the middle line.  This proves the equality as matrix-valued
measures, with every wall included.  None of the steps
inverts \(DF\).
\end{proof}

For a graph \(F(x)=(x,\phi_1(x),\ldots,\phi_r(x))\),
the old coordinate functions have zero Hessian, so
\[
       \mathcal W_{F,B}
           =\sum_{a=1}^r B_{d+a}(F)\,D^2\phi_a.
 \tag{GD5}
\]
The sum is taken before its total variation.  It is the
joint weighted wall measure, not a collection of unrelated
coordinatewise bounds.

If a wall is rational for a lattice \(\Lambda\), let
\(\nu_0\in\Lambda^*\) be its primitive oriented conormal.
When the gradient jump of a scalar coordinate is
\(q\nu_0\), its Hessian, using the ambient lattice measure,
is
\[
                    q\,\nu_0\otimes\nu_0\,\lambda_H,
 \tag{GD6}
\]
where \(\lambda_H\) is the induced lattice measure on the
wall.  Indeed the quotient of the wall and ambient lattice
covolumes is \(|\nu_0|\), which cancels the conversion
from \(\nu_0\) to the Euclidean unit normal in GD3.
Thus the primitive slope \(q\) and the facet carrier are
both retained.  Formula GD6 is a real lattice normalization;
no analytic source multiplicity is replaced by one.

\subsection{The trace version of the strain gap}

The KG estimate extends to noncontinuous branch fields with
specified tangential traces.  More precisely, let
\(g_i\in\mathrm{BD}(B_R^i)\cap L^2(B_R^i)\).  For each pair \(i,j\),
use the bounded trace map on the two half-balls cut out by
\(S_{ij}\).  Suppose that the tangential restrictions of
both one-sided traces of \(g_i\) and both one-sided traces
of \(g_j\) agree almost everywhere on \(D_R^{ij}\).
Then
\[
 R^3|\mathfrak s(c)|\leq
         C\sum_i|Z_i\mathcal L^3-Eg_i|(B_R^i).
 \tag{GD7}
\]
To prove this, apply the same branchwise Korn--Poincare
projection as in KG.  The half-ball estimate is applied
to either chosen one-sided trace.  The trace of the given
continuous \(b_i\) is its usual value, while the assumed
common tangential traces of \(g_i\) cancel in KG7.  The
symmetric-disk argument and the finite stress pairing are
unchanged.  This proof requires no trace on a codimension-two
set and no continuity of the normal components of \(g_i\).

Now let \(F:|\Sigma|\cap\overline{B_R(0)}\to\mathbf R^m\)
be one continuous finite-PL map, with restrictions \(F_i\)
to the five complete source planes of PK.  Let \(B\) be
smooth on a neighborhood of its compact image.  Write \(A_i\) and
\(\mathcal W_i\) for the regular and wall terms of GD3
on \(P_i\).  The piecewise fields \(g_i=F_i^*B\) satisfy
the trace hypothesis of GD7.  Indeed the two restrictions
of \(F\) agree on \(S_{ij}\); their tangential derivatives
there agree almost everywhere, independently of their
one-sided normal derivatives.  Pairing with the continuous
value \(B(F)\) preserves this equality.

\begin{corollary}[A repair must carry a definite wall drift]
\label{cor:r13-graph-repair-drift-lower-bound}
There is a constant \(C>0\) depending only on the PK
configuration such that
\[
 R^3|\mathfrak s(c)|\leq C\sum_i
       \left(\int_{B_R^i}|Z_i-A_i|\,d\mathcal L^3
                          +|\mathcal W_i|(B_R^i)\right).
 \tag{GD8}
\]
If the refined regular strain pulls back exactly to \(Z_i\)
on every old cell, then
\[
        \sum_i|\mathcal W_i|(B_R^i)
                   \geq C^{-1}R^3|\mathfrak s(c)|.
 \tag{GD9}
\]
The same positive lower bound holds for the limit inferior
of the drift masses when the sum of the regular strain
errors tends to zero.
\end{corollary}

\begin{proof}
The pullbacks are locally piecewise smooth, hence belong
to \(\mathrm{BD}\) on the bounded balls, and have the
common tangential traces just checked.  Substitute GD3
into GD7 and use the triangle inequality for measure
variation.  This gives GD8 and its two consequences.
\end{proof}

In particular a graph which is linear on each entire old
source plane has \(\mathcal W_i=0\).  It cannot repair
PK10, even if its ambient space is enlarged.  A piecewise
linear graph can repair it only with nonzero combined
wall drift.  A real subdivision with no change of
coordinate map has zero such drift.

\subsection{An explicit universal real hat-graph repair}

For the particular PK fields there is a concrete graph
repair, before any claim about actual units.  On each
source plane put
\[
                         q(n)=\tfrac12 b_i(n)(n).
 \tag{GD10}
\]
This is one continuous homogeneous piecewise-quadratic
function on \(|\Sigma|\).  On an intersection \(S_{ij}\),
PK9 gives \(b_i(n)(n)=\theta_{ij}(n,n)=0\), so the
values agree between the planes.  On a coefficient cone,
write \(b_i(n)=L_\sigma n\).  Then
\[
                  \operatorname{Hess}(q|_\sigma)
                         =\operatorname{sym}L_\sigma=Z_\sigma.
 \tag{GD11}
\]
This is the regular-cell Hessian; the full Hessian can
also have the wall measures of GD3.

Take a finite common simplicial fan on which \(q\) is
quadratic.  Choose one nonzero rational vector \(v_\rho\)
for each geometric ray, and let \(\phi_\rho\) be its
nodal piecewise-linear coordinate: it takes value one
at \(v_\rho\), value zero at the other selected ray
vectors, and is linear on every cone.  These are
continuous functions on the entire fan.  Define the
ordinary quadratic polynomial
\[
 \begin{split}
 Q(y)&=\sum_\rho q(v_\rho)y_\rho^2\\
 &\quad+\sum_{\{\rho,\rho'\}\text{ spanning a two-face}}
   \bigl(q(v_\rho+v_{\rho'})-q(v_\rho)-q(v_{\rho'})\bigr)
                                      y_\rho y_{\rho'}.
 \end{split}
 \tag{GD12}
\]
Every coefficient is well defined by the common closed-face
values of \(q\).  Polarization in the ray basis of a
simplicial cone proves
\[
                    q=Q((\phi_\rho)_\rho).
 \tag{GD13}
\]
For the graph
\(F(n)=(n,(\phi_\rho(n))_\rho)\), take the smooth ambient
covector \(B=dQ\), with zero old-coordinate components.
It vanishes at the origin and satisfies
\[
                  F^*e(B)|_\sigma=Z_\sigma
 \tag{GD14}
\]
on every old coefficient cone.  It is one ambient covector
on this refined graph, so all of the graph's simultaneous
incident relations hold.  Formula GD9 shows that its old
pullback has unavoidable wall drift when \(\mathfrak s(c)\ne0\).
The graph alone is not asserted to be a balanced actual
source cycle.

\subsection{Actual unit realization on a supplied split SNC core}

There is a finite original-field algebraic construction
which realizes these nodal coordinates on a supplied SNC
core.  Its conclusion deliberately stops at that core.

\begin{lemma}[Simultaneous semilocal boundary generators]
\label{lem:r13-semilocal-boundary-generators}
Let \(V\) be a smooth projective integral variety over the
original finite residue field \(k\), and let
\(D_1,\ldots,D_r\) be distinct prime boundary Cartier
divisors.  Let \(Z\) be any finite set of closed points
meeting every \(D_i\).  There are functions
\(f_1,\ldots,f_r\in k(V)^\times\) and one open
neighborhood of \(Z\) such that
\[
                     \operatorname{div}(f_i)=D_i
                      \quad\hbox{on that neighborhood}.
 \tag{GD15}
\]
In particular
\(\operatorname{ord}_{D_j}(f_i)=\delta_{ij}\) for all
\(i,j\), and every additional zero or pole component
avoids \(Z\).
\end{lemma}

\begin{proof}
A finite set of points of a projective variety lies in an
affine open: use a sufficiently high very ample power,
surjectivity of evaluation at the finite reduced subscheme,
and a section nonzero at each point.  Work in the semilocal
ring \(A=\mathcal O_{V,Z}\) of such an affine open.
Each invertible ideal \(\mathcal O_V(-D_i)\otimes A\)
is free of rank one.  Explicitly, choose a nonzero generator
in each of its finitely many one-dimensional residue
fibers; the Chinese remainder theorem lifts these choices
to one element \(f_i\) of the ideal.  Nakayama's lemma
at every maximal ideal makes this a generator over \(A\).
As the ideal is embedded in the fraction field \(k(V)\),
this is a nonzero rational function.  Its generating
property holds on an open neighborhood of each point of
\(Z\).  Intersect finitely many such neighborhoods for
all \(i\).  Equality of the invertible ideals gives GD15.
Every \(D_j\) meets that neighborhood, so its generic
order is the stated Kronecker delta.  The assertion about
extra components follows from the same divisor equality.
No extension of \(k\) is used.
\end{proof}

\begin{corollary}[Finite actual hat coordinates on the retained core]
\label{cor:r13-actual-hats-on-supplied-snc-core}
Assume the fixed original analytic model germ of R12 TC
is supplied with a split simplicial SNC core on a smooth
projective threefold \(V\) over \(k\).  Suppose its
finitely many retained monomial cones are centered at
triple points of the boundary, with compatible ray labels
\(D_i\), and that its old coordinate map faithfully
parameterizes the geometric core fan.  Use the actual
order vectors of the divisorial valuations as its selected
ray vectors; their old projected vectors need not be
primitive.  Coefficient data on coincident geometric cones
must agree, while all actual source labels and their
calibrations are retained.

There is one finite tuple of actual norm-one units over
the original \(K\), retaining the old coordinates, whose
new moment coordinates, with convention \(-\log|\cdot|\),
equal \(\phi_i\) on every retained
closed old monomial cone, after one common central
shrinking.  Thus any given continuous \(q\) which is quadratic on
each core cone and has polynomial representation GD12 has
the regular-cell identity GD14 on its actual refined
core graph.
\end{corollary}

\begin{proof}
Let \(Z\) contain all retained triple points.  If necessary,
add one closed point on each boundary component not
already meeting \(Z\).  Apply the semilocal lemma.
At a retained triple point with regular normal parameters
\(z_1,z_2,z_3\), the function \(f_i\) is the corresponding
parameter times a unit if \(D_i\) passes through the
point, and is a unit otherwise.  Its order vector is
therefore exactly the desired nodal vector.

Apply the frozen finite same-field residue-unit lifting
result R6 to the finite list \(f_i\).  The TC monomial
comparison evaluates their unit factors with norm one
on all sufficiently near-central positive monomial
parameters.  The finite list of functions and charts
has one common collar and ordinary central shrinking.
Hence their log values are the stated normal parameters.
Continuity extends these equalities to all closed cone
faces, including commensurable radii.  Retaining the old
coordinates retains their old graded rank, so every
old full-source point remains in the new full source.
These points form the asserted refined core graph.

The normal coefficient one is the actual divisorial
order in GD15.  When an old projected ray has a
nontrivial primitive index, that index remains its
projection index; it is not absorbed into a coefficient
or a source weight.  The polynomial identity GD13 and
the ordinary linear chain rule on every core cone now
give GD14.
\end{proof}

\paragraph{The remaining complete-source equation.}
The functions \(f_i\) generally have additional zero and
pole divisors outside \(Z\).  Their new incident branches
and any branches exposed along old double curves remain
in the complete actual refinement.  GD14 proves nothing
about the strain equation on those branches.  To obtain
an actual refinement theorem one must construct or verify
an extension with the correct strains on that entire
new source family, with its labels, weights, primitive
indices and later refinements.  The semilocal lemma does
not supply that extension.  In particular its core graph
cannot replace the absolute cap carrier.  The full
wall drift GD3 and the quantitative lower bound GD9
explain why disregarding those added directions loses
essential data.

\paragraph{Dependencies and scope.}
GD3--GD6 and GD10--GD15 are direct distributional and
finite algebraic calculations.  GD7--GD9 use the reviewed
KG estimate and the same two primary \(\mathrm{BD}\)
inputs.  The conditional analytic core corollary uses
the frozen R6 same-field unit lift and the supplied
R12 TC monomial comparison only on its stated core.
No equality against all original analytic tests or all
GK delta tests is inferred from these real coefficient
identities.
